% !TEX root = ../main.tex
% English ABSTRACT.
\chapter*{Abstract}
\addcontentsline{toc}{chapter}{Abstract}

\begin{center}\bfseries \ThesisTitleEN\end{center}
\noindent\textbf{\StudentEN} \hfill \textbf{Supervisor: \SupervisorEN}\\[10pt]

Motion prediction---forecasting the future trajectories of surrounding vehicles,
pedestrians and cyclists---is a safety-critical component of the autonomous-driving
pipeline that must run in real time on embedded hardware. State-of-the-art
predictors, however, are trained on compute clusters and are too large to run on a
single consumer GPU, placing both ends of the contemporary pipeline out of reach
for an individual researcher. This thesis asks how small a competitive trajectory
predictor can be made before its accuracy degrades, and whether the lost accuracy
can be recovered through \emph{knowledge distillation} without enlarging the model
or worsening the calibration a downstream planner depends on.

The study uses HiVT, a transformer-based Laplace-mixture predictor that is small
enough to be both trained and run on a single GPU, evaluated on the Argoverse~1
benchmark. The accuracy--capacity trade-off is first characterised by sweeping the
embedding width (128, 64, 32, 16) and locating the point at which a from-scratch
student falls measurably below the teacher. The
\emph{mode-permutation problem} is then identified: because HiVT trains its mixture modes with a
winner-takes-all loss, the mode slots of two independently trained models do not
correspond, so any distillation term that aligns modes by index supervises the
student with self-contradictory targets. To resolve this, a
permutation-invariant mixture negative-log-likelihood objective is derived that treats the
teacher's modes as an order-free set of soft targets and supports unequal mode
counts, with a proof of invariance.

Experiments show that a mean-target variant of this objective recovers roughly
$84\%$ of the HiVT-32$\to$HiVT-64 capacity gap ($-9.2\%$ \texttt{minFDE} over a
matched non-distilled baseline) at \emph{zero} added inference cost, but degrades
full-distribution calibration (mixture NLL $+41\%$, calibration error $5\times$) by
discarding the teacher's predictive variance. A distribution-matching objective
that also distils the teacher's per-mode scales removes this penalty entirely,
leaving the student \emph{better calibrated than both the non-distilled baseline
and the teacher} while retaining the full geometric gain. The benefit \emph{grows
as the student shrinks}: at width~16 ($55\times$ smaller than the teacher)
distribution-matching distillation improves \texttt{minFDE} by
$-22.7\%$---roughly $2.5\times$ the width-32 gain---recovering $\sim\!81\%$ of the
width-16$\to$width-32 gap, with calibration improving rather than degrading.
Distillation thus buys close to a full size-class of accuracy for free, and most
where capacity is scarcest.

A final efficiency analysis quantifies the deployment frontier: parameter and
memory savings are fixed and unconditional ($15\times$ at width~32, $55\times$ at
width~16), whereas the single-scene latency speed-up is far sublinear and
batch-dependent (on CPU $\sim\!3\times$ online, rising to $\sim\!5.5\times$ under
modest batching), locating the compression benefit primarily in memory footprint.
Overall, the answer to how small a competitive HiVT can be made is encouraging:
with a permutation-invariant, calibration-preserving distillation loss, a
$55\times$-smaller student reaches roughly the accuracy of an un-distilled model
four times its size at no calibration cost.
